\documentclass[10pt]{article}
\usepackage[margin=0.85in]{geometry}
\usepackage{amsmath,amssymb}
\usepackage[hidelinks]{hyperref}
\setlength{\parindent}{0pt}
\setlength{\parskip}{5pt}
\newcommand{\E}{\mathbb{E}}
\newcommand{\sg}{\operatorname{sg}}
\newcommand{\Cov}{\operatorname{Cov}}
\title{Correcting Tail Bias in Top-\(k\) Score Centering}
\author{}
\date{}

\begin{document}
\maketitle
\vspace{-2.5em}

\paragraph{Abstract.}
Top-$k$ approximations to score centering can introduce systematic update bias when the sampler and trainer disagree on the probability distribution over the unlogged tail. We establish this failure mode through a minimal three-token counterexample and propose an unbiased correction that computes the logged-head contribution exactly while estimating the tail contribution using independent Monte Carlo samples from the sampler's conditional tail distribution. We prove that the corrected estimator recovers full score centering in expectation. A synthetic simulation involving 152,064 tokens further demonstrates that, under sampler-trainer tail mismatch, the top-$k$ approximation can reverse the direction of the expected update relative to full centering. This note presents the problem and the correction, and provides accompanying simulation code and pseudocode for integration with vLLM and SGLang samplers.\footnote{Simulation code and implementation details: \url{https://github.com/mukurgupta/score_center_correcting_bias}.}


\paragraph{Setting.}
We use the notation from Appendix A of Marek and Ryabinin~\cite{marek2026}: $p_v$ and $q_v$ are trainer and sampler probabilities, $s_v=\nabla_\theta\log p_v$, and $y_t\sim q$. All expectations are conditional on a fixed prefix and fixed trainer parameters. Equation numbers below refer to this note unless explicitly identified as equations from the paper.

Exact score centering subtracts $\bar s=\sum_v q_v s_v$. To avoid storing all sampler probabilities, the paper logs probabilities only for the sampler's top-$k$ tokens, represented by the set $H$. For the remaining tokens in the tail set $T$, it estimates the missing probabilities using the trainer. With $Q_T=\sum_{v\in T}q_v$ and $P_T=\sum_{v\in T}p_v$, Eq.~(7) of the paper defines the approximate sampler distribution:
\begin{equation}\label{eq:reconstruction}
 \widehat q_v=\begin{cases}
 q_v,&v\in H,\\
 \rho p_v,&v\in T,
 \end{cases}
 \qquad \rho=Q_T/P_T.
\end{equation}
It assigns the correct total mass to the tail but assumes that the relative probabilities of tail tokens match the trainer's. The implementation therefore subtracts $\widehat{\bar s}=\sum_v\widehat q_v s_v$ instead of $\bar s$ (Eq.~(8) of the paper).


\paragraph{Problem.}
In this note, we target the systematic update error caused by this top-$k$ approximation. Let $\delta$ denote the error in the centering vector. Its effect on the expected update is
\begin{equation}\label{eq:drift}
\begin{aligned}
 \delta&=\bar s-\widehat{\bar s}=\sum_{v\in T}(q_v-\rho p_v)s_v,\\
 \E_q[R(s_{y_t}-\widehat{\bar s})]
 &=\Cov_q(R,s_{y_t})+\E_q[R]\delta.
\end{aligned}
\end{equation}

The extra term $\E_q[R]\delta$ appears because of the top-$k$ approximation. If it is nonzero, the approximation introduces an additional update beyond the reward--score covariance. We first show a case where the error is nonzero, then propose a simple fix that removes it in expectation.


\paragraph{Counterexample (where $\delta$ matters).}
Consider a toy vocabulary with three tokens $(A,B,C)$. Let $p=(0.990,0.009,0.001)$, $q=(0.990,0.001,0.009)$, and let the top-$k$ head be $H=\{A\}$. The trainable parameters are the three softmax logits, so $s_v=e_v-p$, where $e_v$ is the token's unit vector. The total tail mass (for tokens $B$ and $C$) is equal for the sampler and trainer, implying $\rho=1$, $\widehat q=p$, and $\widehat{\bar s}=0$.

As in Section 4 of the paper, give every token reward $R=1$. The reward factor then drops out of the update, and the expected-reward objective is constant with zero gradient. Nevertheless, subtracting $\widehat{\bar s}$ produces a nonzero expected update:
\begin{equation}\label{eq:toy}
 \E_q[s_{y_t}-\widehat{\bar s}]=q-p=(0,-0.008,0.008),
 \qquad \E_q[s_{y_t}-\bar s]=0.
\end{equation}

See \hyperref[app:toy]{appendix proof~1} for the derivation of \eqref{eq:toy}. This happens because the probabilities of tokens $B$ and $C$ differ between $p$ and $q$, and both tokens are outside the top-$k$ set $H$. Thus 99\% head coverage does not guarantee a zero mean update for constant rewards.

\paragraph{Proposed fix.}
We estimate the missing tail contribution using a few tokens sampled from the sampler's conditional tail distribution. Before discarding the full sampler distribution, draw $m\ge1$ auxiliary token IDs $\mathbf z=(z_1,\ldots,z_m)$ from $q$ restricted to the tail set $T$. Store these IDs rather than their sampler probabilities. For the Monte Carlo estimate, draw independently with replacement, and independently of the actual rollout and reward, conditional on the prefix.

We define a Monte Carlo centering term $\bar s_{\rm MC}$ to use in place of $\widehat{\bar s}$:

\begin{equation}\label{eq:estimator}
 \Pr(z_j=v)=\frac{q_v}{Q_T}\quad(v\in T),\qquad
 \bar s_{\rm MC}=\sum_{v\in H}q_v s_v+\frac{Q_T}{m}\sum_{j=1}^m s_{z_j}.
\end{equation}
The derivation of \eqref{eq:estimator} is given in \hyperref[app:unbiased]{appendix proof~2}. If $Q_T=0$, the head is already exact and no tail draws are needed. The estimate and the resulting update satisfy
\begin{equation}\label{eq:guarantee}
 \E_{\mathbf z}[\bar s_{\rm MC}]=\bar s,\qquad
 \E_{y_t,R,\mathbf z}[R(s_{y_t}-\bar s_{\rm MC})]=\Cov_q(R,s_{y_t}).
\end{equation}
The first claim in \eqref{eq:guarantee} is proved in \hyperref[app:unbiased]{appendix proof~2}; independence gives $\E[R(\bar s_{\rm MC}-\bar s)]=\E[R]\E_{\mathbf z}[\bar s_{\rm MC}-\bar s]=0$, proving the second. The estimator applies Liu et al.'s existing construction: sum the head exactly and sample the tail~\cite{liu2019}.


\paragraph{Implementation.}
The corresponding scalar loss, modifying Eq.~(6) of the paper, is
\begin{equation}\label{eq:loss}
 L=-\sg(R)\left[\log p_{y_t}-\sum_{v\in H}\sg(q_v)\log p_v
 -\frac{\sg(Q_T)}{m}\sum_{j=1}^m\log p_{z_j}\right].
\end{equation}
Store $m$ extra token IDs and evaluate their trainer log-probabilities. 
In \texttt{README.md}, we include a pseudo-code of the minimal changes required in the sampler implemented using vLLM or SGLang. The variance cost is derived in \hyperref[app:variance]{Appendix~3}.


\paragraph{Experiments.}
We simulate this toy example using \texttt{simulate\_toy.py}, which enumerates the possible draws to compute exact mean updates and checks them with 500,000 Monte Carlo trials. The exact mean update is $(0,-0.008,0.008)$ for approximate top-$k$ centering and zero for both full centering (the exact reference) and the proposed Monte Carlo fix.

While the toy example is interpretable, it is far from real-world scenarios. We also include a simulation with a vocabulary of 152,064 tokens showing a case closer to real-world settings.

\paragraph{Practical Implications.}
The hypothesis under consideration in this note is a systematic sampler--trainer difference across tail tokens. 

Although speculative, possible scenarios in which this effect may have a major impact include:
\begin{itemize}
    \item Multilingual or specialist text: related script tokens or technical terms may exhibit coherent tail mismatches.
    \item Code and tool arguments: tail mismatches may bias updates for rare identifiers or argument values.
\end{itemize}


\paragraph{AI Assistance.}
We used AI assistants to polish the writing of this note and to understand the literature. After identifying the top-K limitation addressed in this note, we used AI tools to craft proofs and synthetic simulations, which we thoroughly verified.







\clearpage
\section*{Appendix: Proofs}
\phantomsection\label{app:proofs}
\setcounter{equation}{0}
\renewcommand{\theequation}{A.\arabic{equation}}

The proofs below establish the counterexample, unbiasedness of the estimated centering term, and its variance cost. All expectations are conditional on the same fixed prefix and fixed trainer parameters.

\subsection*{1. The constant-reward counterexample}
\phantomsection\label{app:toy}
We prove \eqref{eq:toy}, the difference between the expected updates under approximate and full centering for constant rewards.

For directly trainable softmax logits ($\ell_v$):
\begin{equation}
\begin{aligned}
 p_v&=\frac{\exp(\ell_v)}{\sum_u\exp(\ell_u)},\\
 \frac{\partial\log p_v}{\partial\ell_u}
 &=\mathbf 1\{u=v\}-p_u,
 \qquad s_v=\nabla_\ell\log p_v=e_v-p.
\end{aligned}
\end{equation}
$e_v$ is the unit vector for token $v$. In the counterexample, equal tail masses give $\rho=1$ and $\widehat q=p$. Consequently,
\begin{equation}
\begin{aligned}
 \bar s&=\sum_vq_v(e_v-p)=q-p=(0,-0.008,0.008),\\
 \widehat{\bar s}&=\sum_vp_v(e_v-p)=0.
\end{aligned}
\end{equation}
With reward one for every token, the true objective has zero gradient:
\begin{equation}
 J(\ell)=\sum_vp_v\cdot 1=1,\qquad \nabla_\ell J=0.
\end{equation}
But the two centering rules give different expected updates:
\begin{equation}
\begin{aligned}
 \E_q[s_{y_t}-\widehat{\bar s}]
 &=\bar s-\widehat{\bar s}=(0,-0.008,0.008),\\
 \E_q[s_{y_t}-\bar s]&=0.
\end{aligned}
\end{equation}
This proves \eqref{eq:toy}.

\subsection*{2. Unbiased estimation of the tail contribution}
\phantomsection\label{app:unbiased}
We derive the estimator in \eqref{eq:estimator} and prove the first claim in \eqref{eq:guarantee}: its expected centering vector equals $\bar s$. The missing contribution from the tail set can be written as
\begin{equation}\label{eq:tail-decomposition}
 \sum_{v\in T}q_vs_v
 =Q_T\sum_{v\in T}\frac{q_v}{Q_T}s_v
 =Q_T\E_{z\sim q(\cdot\mid T)}[s_z].
\end{equation}
Replacing this conditional expectation by a sample average gives \eqref{eq:estimator}. Each draw has the conditional mean in \eqref{eq:tail-decomposition}. By linearity of expectation,
\begin{equation}\label{eq:unbiased-proof}
\begin{aligned}
 \E_{\mathbf z}[\bar s_{\rm MC}]
 &=\sum_{v\in H}q_vs_v+\frac{Q_T}{m}\sum_{j=1}^m\E[s_{z_j}]\\
 &=\sum_{v\in H}q_vs_v+\frac{Q_T}{m}\,m\sum_{v\in T}\frac{q_v}{Q_T}s_v\\
 &=\sum_vq_vs_v=\bar s.
\end{aligned}
\end{equation}
This proves the first claim in \eqref{eq:guarantee}.

\subsection*{3. Variance cost and the benefit of summing the head}
\phantomsection\label{app:variance}
Assume finite second moments. For a random vector $X$, write
$\Cov(X)=\E[(X-\E X)(X-\E X)^\top]$ for its covariance matrix.
The head contribution is fixed at the prefix. Independence of the $m$ tail draws therefore gives
\begin{equation}
\begin{aligned}
 \Cov_{\mathbf z}(\bar s_{\rm MC})
 &=\frac{Q_T^2}{m^2}\sum_{j=1}^m\Cov(s_{z_j})\\
 &=\frac{Q_T^2}{m}\Cov_{q(\cdot\mid T)}(s_z),
\end{aligned}
\end{equation}
If $Q_T=0$, the head is exact and this covariance is zero.

An alternative is to draw $x_1,\ldots,x_m$ independently from the full sampler distribution $q$ and average their scores:
\begin{equation}
 \bar s_q=\frac1m\sum_{j=1}^m s_{x_j},\qquad
 \E[\bar s_q]=\bar s,\qquad
 \Cov(\bar s_q)=\frac1m\Cov_q(s_x).
\end{equation}
This is also unbiased, provided the draws are independent of the rollout and reward. Summing the head exactly instead leaves only conditional-tail uncertainty, multiplied by $Q_T^2$. This factor is desirable when the head contains most of the probability mass; it is not generally the exact variance ratio between the two estimators, since their score covariances differ.

Finally, define $g_{\rm full}=R(s_{y_t}-\bar s)$ and $g_{\rm MC}=R(s_{y_t}-\bar s_{\rm MC})$. With $\delta_{\rm MC}=\bar s_{\rm MC}-\bar s$, we have $g_{\rm MC}=g_{\rm full}-R\delta_{\rm MC}$. The error $\delta_{\rm MC}$ has zero mean and is independent of $(y_t,R)$, so the cross-covariances vanish. Hence
\begin{equation}\label{eq:update-variance}
 \Cov(g_{\rm MC})=\Cov(g_{\rm full})
 +\E[R^2]\frac{Q_T^2}{m}\Cov_{q(\cdot\mid T)}(s_z).
\end{equation}
This is an added variance cost relative to full centering. The sampled estimator need not have lower total error than the deterministic top-$k$ approximation; when $p=q$, that approximation is already exact while tail sampling can still add noise.

\begin{thebibliography}{9}
\setlength{\itemsep}{0pt}
\bibitem{marek2026} M. Marek and M. Ryabinin. \emph{Score Centering Stabilizes Off-policy Reinforcement Learning}. 2026. \url{https://arxiv.org/abs/2609.20807}.
\bibitem{liu2019} R. Liu et al. \emph{Rao-Blackwellized Stochastic Gradients for Discrete Distributions}. ICML, 2019. \url{https://proceedings.mlr.press/v97/liu19c.html}.
\end{thebibliography}
\end{document}
